\subsection{网格图：一个 lane 的八条 A 行、四条 B 列}
继续看 Warp 0 的 lane 11：laneRow=1、laneCol=3。为放大格子，下图只展示 Warp 0 对应的 A 前 32 行、B 前 32 列以及 C 的左上 32×32 区域；\textbf{它们仍来自全 Block 的 Shared tile，并非 Warp 私有数组。}
\begin{figure}[H]\centering
\begin{tikzpicture}[x=1cm,y=1cm]
\GemmReadGrid{32}{32}{1,5,9,13,17,21,25,29}{3,11,19,27}{6}
\node[lbl,text width=4.7cm] at (2.2,5) {Shared tileA：前 32 行\\32×16};
\node[lbl,text width=4.7cm] at (8,5) {Shared tileB：前 32 列\\16×32};
\node[lbl,text width=4.7cm] at (13.8,5) {Warp 0 的 C：32×32\\本地行、列 0--31};
\node[lbl,text width=4.7cm] at (2.2,-1) {固定八行：1、5、9、13、\\17、21、25、29\\k 沿各行从 0 走到 15};
\node[lbl,text width=4.7cm] at (8,-1) {固定四列：3、11、19、27\\k 沿各列从 0 走到 15};
\node[lbl,text width=4.7cm] at (13.8,-1) {只更新 8 行×4 列的交叉点\\共 32 个分散输出};
\node[lbl,text width=16.2cm] at (8,-2.3) {V6：Warp 0 的 lane 11，laneRow=1、laneCol=3，深橙格示例为 k=6。\\每个 k：8 个 A、4 个 B 交叉配对，给 32 个 C 各增加一项。};
\end{tikzpicture}

\caption{与 V1/V2 相同，A 沿行遍历 k、B 沿列遍历 k。区别是 lane 11 同时负责八行与四列的 32 个交叉输出，而不是一块连续的 8×4 区域。网格每格一个元素，横纵缩放随各自维度变化。}
\end{figure}
这个 lane 在每个 k 执行
\[
accum[i][j]\mathrel{+}=tileA[1+4i][k],tileB[k][3+8j],\quad i=0\ldots7,\ j=0\ldots3.
\]
例如 $accum[2][2]$ 对应本地 $C[9,19]$。k=6 时加 $tileA[9][6]\times tileB[6][19]$；本轮完整贡献为
\[
\Delta C[9,19]=\sum_{k=0}^{15}tileA[9][k],tileB[k][19].
\]
若 Block 起点仍取 $(64,128)$，这个输出就是全局 $C[73,147]$。剩余 31 个橙色输出同样各完成 16 项点积，其他格子由 Warp 0 的其他 lane 负责。
\takeaway{V2 每 k 用 4A+4B 更新 16 个输出；V6 每 k 用 8A+4B 更新 32 个输出。K 循环仍是 16 步，没有少算点积，也不是一次读出整行整列。}
\pagebreakhere
